% ============================================================
% DS MOD08 : loi binomiale, loi de Poisson, dénombrement
% BTS ATI / CRSA 1re année -- durée 1 h
% Pôle Formation UIMM - CVDL - S. JAUBERT
% Corrigé et barème : MOD08_DS_Binomiale_Poisson_ATI_CRSA_Corrige_v1_0
% ============================================================
\documentclass[a4paper,11pt]{article}
\usepackage[utf8]{inputenc}
\usepackage[T1]{fontenc}
% babel-french absent sur certains postes : chargement conditionnel
\IfFileExists{french.ldf}{\usepackage[french]{babel}}{}
\usepackage{helvet}
\usepackage{courier}
\renewcommand{\familydefault}{\sfdefault}
\usepackage[top=3cm, bottom=2.5cm, left=2cm, right=2cm, headheight=2cm]{geometry}
\usepackage{amsmath,amssymb,amsfonts}
\usepackage{graphicx}
\usepackage[dvipsnames,table]{xcolor}
\usepackage{tikz}
\usepackage[most]{tcolorbox}
\usepackage{fancyhdr}
\usepackage{enumitem}
\usepackage{array,needspace}
\usepackage{titlesec}
\usepackage{hyperref}

\definecolor{uimmblue}{HTML}{005C85}
\definecolor{uimmred}{HTML}{D31326}
\definecolor{uimmlightblue}{HTML}{E8F1F8}
\definecolor{uimmgray}{HTML}{333333}
\definecolor{uimmlightgray}{HTML}{F5F5F5}
\hypersetup{colorlinks=true, linkcolor=uimmblue, urlcolor=uimmblue}

\titleformat{\section}{\Large\bfseries\color{uimmblue}}{}{0pt}{}[\vspace{2pt}{\color{uimmblue}\titlerule[1.2pt]}]
\setlist[enumerate,1]{label={\color{uimmblue}\textbf{\arabic*.}}, leftmargin=*}
\setlist[enumerate,2]{label=\alph*), leftmargin=*}
\setlist[itemize,1]{label={\color{uimmblue}\textbullet}}

\newtcolorbox{exo}[1]{colback=white, colframe=uimmblue, fonttitle=\bfseries\color{white},
  coltitle=white, title={#1}, boxrule=1pt, arc=2pt,
  left=8pt, right=8pt, top=5pt, bottom=6pt, before skip=10pt, after skip=6pt}
\newtcolorbox{consigne}{colback=uimmlightgray, colframe=uimmgray, boxrule=0.8pt, arc=3pt,
  left=8pt, right=8pt, top=6pt, bottom=6pt}

\newcommand{\pts}[1]{\hfill\mbox{\color{uimmred}\small\textbf{(#1)}}}
\newcommand{\calc}{\,{\color{uimmred}\small\textbf{[calculatrice]}}}
% Zone de réponse lignée : #1 = nombre de lignes
\newcommand{\zone}[1]{\par\smallskip\noindent\begin{tikzpicture}
  \foreach \i in {1,...,#1}{\draw[uimmgray!30] (0,-0.85*\i) -- (\linewidth-1pt,-0.85*\i);}
  \end{tikzpicture}\par\medskip}
% Cadre vide (arbre pondéré) : #1 = hauteur
\newcommand{\cadre}[1]{\par\smallskip\noindent\begin{tikzpicture}
  \draw[uimmgray!40, rounded corners=3pt] (0,0) rectangle (\linewidth-1pt,-#1);
  \end{tikzpicture}\par\medskip}
\newcommand{\EnTete}[2]{%
\pagestyle{fancy}\fancyhf{}\renewcommand{\headrulewidth}{0pt}
\fancyhead[L]{\raisebox{-0.5\height}{\includegraphics[height=1.5cm]{logo_uimm_placeholder.jpg}}}
\fancyhead[C]{\begin{minipage}{8cm}\centering
  {\small\color{uimmblue}\textbf{Pôle Formation UIMM - CVDL}}\\
  {\footnotesize\color{uimmgray} BTS ATI / CRSA 1\textsuperscript{re} année\\ Mathématiques}\end{minipage}}
\fancyhead[R]{{\small\color{uimmred}\textbf{MOD08 -- #1}}}
\fancyfoot[C]{{\color{uimmblue}\rule{\textwidth}{1pt}}\\[2pt]
  {\footnotesize\color{uimmgray} Pôle Formation UIMM -- Centre-Val de Loire \hfill S. JAUBERT \hfill #2 -- \thepage}}}
\EnTete{DS Binomiale -- Poisson}{DS 1 h}
\newtcolorbox{rappel}{colback=uimmlightblue, colframe=uimmblue, fonttitle=\bfseries\color{white},
  coltitle=white, title={Rappels de cours}, boxrule=1pt, arc=2pt,
  left=8pt, right=8pt, top=5pt, bottom=6pt, before skip=8pt, after skip=8pt}

\begin{document}

\begin{tcolorbox}[colback=uimmred, colframe=uimmred, sharp corners, boxrule=0pt,
  left=14pt, right=14pt, top=10pt, bottom=10pt]
{\color{white}\LARGE\bfseries Devoir surveillé -- MOD08}\\[4pt]
{\color{white}\large Loi binomiale, loi de Poisson, dénombrement \hfill Durée~: 1 heure}
\end{tcolorbox}

\vspace{2pt}
\noindent{\small\begin{tabular}{@{}>{\bfseries\color{uimmblue}}l@{\hspace{8pt}}p{13.8cm}@{}}
Nom, prénom~: & \dotfill \hspace{1cm}{\bfseries\color{uimmblue}Note~:}\ \dots\dots\ / 20 \\
\end{tabular}}

\begin{consigne}
\textbf{Consignes.} Calculatrice autorisée en mode examen, aucun document. Répondre sur le sujet, dans les zones prévues.
\begin{itemize}[leftmargin=12pt, itemsep=1pt, topsep=2pt]
  \item Toute loi utilisée doit être \textbf{nommée avec ses paramètres}.
  \item Pour chaque probabilité demandée~: écrire l'évènement à l'aide de la variable aléatoire, \textbf{puis} le résultat. Un résultat seul ne rapporte pas tous les points.
  \item Probabilités arrondies à $10^{-4}$. Le barème est donné à titre indicatif.
\end{itemize}
\end{consigne}

\begin{rappel}
\textbf{Coefficient binomial.} Nombre de façons de choisir $k$ objets parmi $n$, sans tenir compte de l'ordre~:
\[
\binom{n}{k} = \frac{n!}{k!\,(n-k)!}.
\]
\textbf{Loi binomiale $\mathcal{B}(n\,;p)$.} On répète $n$ fois, de façon identique et indépendante, une épreuve à deux issues (succès de probabilité $p$, échec de probabilité $1-p$). La variable $X$ qui compte le nombre de succès suit la loi $\mathcal{B}(n\,;p)$. Elle prend les valeurs $0, 1, \dots, n$.
\[
P(X = k) = \binom{n}{k}\,p^{k}\,(1-p)^{n-k},\qquad E(X) = np,\qquad \sigma(X) = \sqrt{np(1-p)}.
\]
\textbf{Loi de Poisson $\mathcal{P}(\lambda)$.} Elle modélise un nombre d'évènements rares (pannes, défauts, arrivées) sur une durée ou une zone donnée. Elle prend les valeurs $0, 1, 2, \dots$
\[
P(X = k) = e^{-\lambda}\,\frac{\lambda^{k}}{k!},\qquad E(X) = \lambda,\qquad \sigma(X) = \sqrt{\lambda}.
\]
\textbf{Approximation.} Si $n \geq 30$, $p \leq 0{,}1$ et $np \leq 10$, la loi $\mathcal{B}(n\,;p)$ peut être approchée par la loi $\mathcal{P}(\lambda)$ avec $\lambda = np$.

\smallskip
\textbf{Calculatrice.} La commande «~Fdp~» (ou «~pd~») donne $P(X = k)$. La commande «~FRép~» (ou «~cd~») donne $P(X \leq k)$.\\
{\small TI~: \texttt{binomFdp(n,p,k)}, \texttt{binomFRép(n,p,k)}, \texttt{poissonFdp($\lambda$,k)}, \texttt{poissonFRép($\lambda$,k)}. \quad Casio~: \texttt{Bpd}, \texttt{Bcd}, \texttt{Ppd}, \texttt{Pcd}.}
\end{rappel}

% ------------------------------------------------------------
\newpage
\section{Exercice 1 -- Test de cartes d'entrées/sorties \hfill {\normalsize\color{uimmred}4 points}}
\begin{exo}{Contexte}
L'armoire de commande d'une cellule robotisée contient $15$ cartes d'entrées/sorties~: $5$ cartes \textbf{de sécurité} et $10$ cartes \textbf{standard}. Pour un test de maintenance préventive, l'automaticien tire au sort $4$ cartes différentes. Tous les groupes de $4$ cartes ont la même probabilité d'être tirés.
\end{exo}
\begin{enumerate}
  \needspace{4.8cm}
  \item Combien de groupes de $4$ cartes peut-on former~? \pts{0,75 pt}
  \zone{3}
  \needspace{5.6cm}
  \item Combien de groupes contiennent au plus une carte de sécurité~? \pts{1 pt}
  \zone{4}
  \needspace{5.6cm}
  \item Calculer la probabilité que le groupe tiré contienne au moins une carte de sécurité. \pts{1 pt}
  \zone{4}
  \needspace{5.6cm}
  \item Un collègue affirme~: «~Le contraire de \textit{le groupe contient au moins deux cartes de sécurité} est \textit{le groupe contient au plus deux cartes de sécurité}.~» A-t-il raison~? Justifier, puis donner le contraire correct. \pts{0,75 pt}
  \zone{4}
  \needspace{4.8cm}
  \item Les $4$ cartes tirées sont testées l'une après l'autre sur le banc de test. Combien de séquences de test différentes (choix des cartes et ordre de passage) peut-on établir à partir des $15$ cartes de l'armoire~? \pts{0,5 pt}
  \zone{3}
\end{enumerate}

% ------------------------------------------------------------
\newpage
\section{Exercice 2 -- Cartes électroniques \hfill {\normalsize\color{uimmred}8 points}}
\begin{exo}{Contexte}
En sortie de ligne, un banc automatique contrôle des cartes électroniques de commande. On sait que $3\,\%$ des cartes produites présentent un défaut. On prélève $40$ cartes dans la production. La production est assez importante pour assimiler ce prélèvement à un tirage avec remise.\\
On note $X$ le nombre de cartes \textbf{sans défaut} dans le prélèvement.
\end{exo}
\begin{enumerate}
  \needspace{6.5cm}
  \item Justifier que $X$ suit une loi binomiale. Préciser ses paramètres. \pts{1,5 pt}
  \zone{5}
  \needspace{5.6cm}
  \item Calculer l'espérance $E(X)$ et l'écart-type $\sigma(X)$. Interpréter $E(X)$ dans le contexte. \pts{1,5 pt}
  \zone{4}
  \needspace{4cm}
  \item Les cartes sont numérotées dans l'ordre du prélèvement, de $1$ à $40$.
  \begin{enumerate}
    \needspace{3.2cm}
    \item De combien de façons peut-on placer $39$ cartes sans défaut parmi les $40$ cartes du prélèvement~? \pts{0,5 pt}
    \zone{2}
    \needspace{4.8cm}
    \item En déduire la probabilité que le prélèvement contienne exactement $39$ cartes sans défaut. Détailler le calcul \textbf{sans} utiliser les fonctions de loi de la calculatrice. \pts{1 pt}
    \zone{3}
  \end{enumerate}
  \needspace{4cm}
  \item Pour chacun des évènements suivants, traduire l'évènement à l'aide de $X$, puis calculer sa probabilité.
  \begin{enumerate}
    \needspace{3.2cm}
    \item Aucune carte du prélèvement ne présente de défaut. \pts{0,5 pt}
    \zone{2}
    \needspace{4.8cm}
    \item Le prélèvement contient au plus $2$ cartes défectueuses. \pts{1 pt}
    \zone{3}
    \needspace{4.8cm}
    \item Le prélèvement contient plus de $3$ cartes défectueuses. \pts{1 pt}
    \zone{3}
    \needspace{4.8cm}
    \item Le prélèvement contient au moins une carte défectueuse, mais pas plus de $3$. \pts{1 pt}
    \zone{3}
  \end{enumerate}
\end{enumerate}

% ------------------------------------------------------------
\needspace{12cm}
\section{Exercice 3 -- Alarmes d'un robot de palettisation \hfill {\normalsize\color{uimmred}6 points}}
\begin{exo}{Contexte}
L'automate qui pilote un robot de palettisation enregistre les alarmes de défaut (bourrage, détection de présence, défaut de préhension). Ces alarmes surviennent de façon aléatoire et indépendante. Le journal de l'automate montre qu'il se produit \textbf{en moyenne $3{,}2$ alarmes par poste de travail de $8$ heures}. On admet que le nombre d'alarmes au cours d'un poste suit une loi de Poisson.
\end{exo}
\begin{enumerate}
  \needspace{4.8cm}
  \item Définir une variable aléatoire adaptée à la situation. Donner sa loi et justifier la valeur de son paramètre. \pts{1 pt}
  \zone{3}
  \needspace{3.2cm}
  \item Calculer l'écart-type de cette variable aléatoire. \pts{0,5 pt}
  \zone{2}
  \needspace{4.8cm}
  \item Calculer la probabilité qu'un poste compte exactement $2$ alarmes. Détailler le calcul à l'aide de la formule. \pts{1 pt}
  \zone{3}
  \needspace{4.8cm}
  \item Calculer la probabilité qu'un poste compte plus de $4$ alarmes. \pts{1 pt}
  \zone{3}
  \needspace{7.4cm}
  \item Le technicien juge un poste «~normal~» lorsque le nombre d'alarmes reste dans l'intervalle $[E - \sigma\,;E + \sigma]$, où $E$ est l'espérance et $\sigma$ l'écart-type de la variable aléatoire. Donner les nombres d'alarmes possibles pour un poste normal, puis calculer la probabilité qu'un poste soit normal. \pts{1,5 pt}
  \zone{6}
  \needspace{5.6cm}
  \item On observe le robot pendant $5$ postes. On admet que les nombres d'alarmes d'un poste à l'autre sont indépendants. Calculer la probabilité qu'aucun de ces $5$ postes ne compte plus de $4$ alarmes. \pts{1 pt}
  \zone{4}
\end{enumerate}

% ------------------------------------------------------------
\needspace{10cm}
\section{Exercice 4 -- Contrôle de connecteurs \hfill {\normalsize\color{uimmred}2 points}}
\begin{exo}{Contexte}
Une machine de sertissage produit des connecteurs en grande série. $0{,}4\,\%$ des connecteurs produits sont défectueux. On prélève $750$ connecteurs dans la production. Le prélèvement est assimilé à un tirage avec remise. On s'intéresse au nombre de connecteurs défectueux dans le prélèvement.
\end{exo}
\begin{enumerate}
  \needspace{4.8cm}
  \item Donner la loi exacte du nombre de connecteurs défectueux. Justifier qu'on peut l'approcher par une loi de Poisson, dont on précisera le paramètre. \pts{1 pt}
  \zone{3}
  \needspace{4.8cm}
  \item À l'aide de cette loi de Poisson, calculer la probabilité que le prélèvement contienne moins de $2$ connecteurs défectueux. \pts{1 pt}
  \zone{3}
\end{enumerate}

\end{document}
